\section{Modèle spectral des zéros}
\label{sec:modele_spectral}

Cette section présente l'armature analytique et opératorielle. Quatre
résultats principaux~: expression close des amplitudes $(\kappap)$,
représentation Fredholm exacte de $\Rres$, opérateur canalisé pour le
couplage intra-cascade, et formule de trace explicite.

\subsection{Expression close des amplitudes $\kappap(s)$}
\label{sec:kappa_p}

\begin{definition}
\label{def:kappa_p}
Pour chaque premier $p$ et chaque $s \in \Cplx$, on pose
\begin{equation}
\boxed{\;\kappap(s) \;=\; 1 \,-\, \big(1 - p^{-s}\big)\,\exp\!\big(\Ap\,p^{-s}\big)\;}
\label{eq:kappa_p}
\end{equation}
\end{definition}

\noindent Par construction,
\[
1 \,-\, \kappap(s) \;=\; \big(1 - p^{-s}\big)\,\exp\!\big(\Ap\,p^{-s}\big),
\]
et le développement asymptotique pour $p \to \infty$ donne
\begin{equation}
\kappap(s) \;=\; (1 - \Ap)\,p^{-s} \,+\, O\!\big(p^{-2s}\big).
\label{eq:kappa_asymp}
\end{equation}

\paragraph{Interprétation.}
Le scalaire fondamental de la cascade n'est ni $p^{-s}$ (terme d'Euler)
ni $\exp(\Ap p^{-s})$ (compensation bifurquée), mais leur écart $\kappap$.
On a la lecture
\[
1 \,-\, \kappap \;=\; \text{survie résiduelle},
\qquad
\kappap \;=\; \text{amplitude spectrale du résidu}.
\]

\subsection{Représentation Fredholm exacte}
\label{sec:fredholm}

Soit $\DRop(s)$ l'opérateur diagonal sur $\ell^{2}(\Primes)$ ($\Primes$
l'ensemble des premiers) défini par $\DRop(s)\,e_p = \kappap(s)\,e_p$.

\begin{theoreme}[Représentation Fredholm de $\Rres$]
\label{thm:fredholm}
Pour $\Reop(s) > 1$, l'opérateur $\DRop(s)$ est de classe trace, et
\begin{equation}
\boxed{\;\det\!\big(I - \DRop(s)\big)^{-1} \;=\; \Rres(s)\;.}
\label{eq:fredholm}
\end{equation}
\end{theoreme}

\begin{proof}
Pour $\Reop(s) > 1$, $\sum_p |\kappap(s)| < \infty$ par le développement
asymptotique~\eqref{eq:kappa_asymp} ; l'opérateur diagonal est donc de
classe trace. Le calcul direct donne
\begin{align*}
\det\!\big(I - \DRop(s)\big)^{-1}
&= \prod_p (1 - \kappap)^{-1}
= \prod_p (1 - p^{-s})^{-1}\,\exp(-\Ap\,p^{-s})\\
&= \zeta(s)\,\exp\!\Big(-\sum_p \Ap\,p^{-s}\Big)
= \Rres(s),
\end{align*}
en utilisant le produit eulérien classique pour $\zeta$ et la définition
de $\Rres$ en~\eqref{eq:Rres_def}.
\end{proof}

\paragraph{Portée de la représentation.}
La preuve est algébrique directe~: \eqref{eq:fredholm} est une
\emph{réécriture déterminantielle} de la définition de $\Rres$. Sa valeur
est de fournir la \emph{cible spectrale} que tout opérateur PT canonique
candidat doit reproduire ; ce n'est pas un argument indépendant pour la
structure spectrale de $\Rres$.

\paragraph{Vérification numérique.}
Pour $s = 1.3 + 14.134725\,i$, $P = 10^{6}$ premiers tronqués, l'erreur
résiduelle entre $\Rres(s)$ calculée par le produit Fredholm
\eqref{eq:fredholm} et $\Rres(s)$ calculée à partir de $\zeta$ par
\texttt{mpmath} est de l'ordre de $2.0 \cdot 10^{-12}$. À $\sigma = 1.1$
la précision tombe à $\sim 1.3 \cdot 10^{-3}$, ce qui reflète l'approche
de la frontière de convergence absolue, non un défaut de la
représentation.

\begin{remarque}[Hypothèse de non-annulation]
\label{rem:non_annulation}
L'identification $\mathrm{z\acute{e}ros}(\Rres) = \mathrm{z\acute{e}ros}(\zeta)$
utilisée dans toute la suite suppose
$\zetaplus(s)\,\zetaminus(s) \neq 0$ sur les zones d'intérêt (en
particulier la ligne critique). Cette hypothèse est vérifiée
numériquement sur les zones testées mais n'est pas démontrée ; elle est
\textbf{conjecturale} et doit être traitée comme telle
(cf.~\S\ref{sec:limites_ouvertures}).
\end{remarque}

\subsection{Opérateur canalisé pour le couplage $\sum_p \Ap p^{-s}$}
\label{sec:operateur_canalise}

Pour modéliser le couplage intra-cascade (la somme $\sum_p \Ap p^{-s}$)
comme spectre d'un opérateur canonique \emph{sans paramètre ajustable},
on construit
\begin{equation}
\Hilb_{\mathrm{PT}}
\;=\;
L^{2}(\Sigma_{\mathrm{pers}}, du) \,\otimes\, L^{2}\!\big(\widehat{\Integers/M\Integers}\big).
\label{eq:H_PT_def}
\end{equation}
La première composante est la coordonnée de Mellin continue sur
$\Sigma_{\mathrm{pers}}$ (surface de persistance, de genre $14$ dans le
corpus PT \cite{PT_GeoSpectral}) ; la seconde est discrète, sur les
caractères du groupe multiplicatif modulo $M$ (cascade par théorème
chinois des restes, $M$ produit primorial). Le couplage inter-fenêtre est
construit par noyau de Bergman régularisé
\begin{equation}
B_{\mathrm{BF}}(z_i, z_j)
\;=\;
\frac{\eps_i\,\eps_j}{|z_i - z_j|^{2} + \eps_i\,\eps_j},
\qquad
\eps_i \;=\; \frac{\epsstar}{\rho_i}.
\label{eq:Bergman_BF}
\end{equation}

\begin{lemme}[Détermination géométrique de $\epsstar$]
\label{lem:epsstar_geom}
La covariance de Bergman régularisée à résolution maximale
($\eps_{\mathrm{reg}} = 10^{-3}$, paquets de taille premier individuel)
sur la mesure atomique $\pi_{\log}$ de la coordonnée de Mellin donne
\begin{equation}
\epsstar^{\mathrm{geom}} \;=\; 1{,}51.
\label{eq:epsstar_geom}
\end{equation}
\end{lemme}

\begin{proof}[Esquisse]
La covariance Bergman régularisée du noyau~\eqref{eq:Bergman_BF} sur la
mesure atomique $\pi_{\log} = \sum_p \delta(u - \log p)$ est, à la
résolution $\eps_{\mathrm{reg}}$, la moyenne pondérée
$\mathrm{Cov}(\eps_{\mathrm{reg}}) = \int\!\!\int B_{\mathrm{BF}}(u, u')^{2}\,
\pi_{\log}(du)\,\pi_{\log}(du')$. Le sweep numérique sur les valeurs de
$\eps_{\mathrm{reg}}$ par paliers $10^{-3}, 10^{-2}, 10^{-1}, 1$ donne
une covariance respectivement de $0{,}0252$, $0{,}0794$, $0{,}2408$ et
$0{,}680$, soit un facteur $\epsstar^{\mathrm{geom}} = 1{,}51$ à la
résolution maximale (référence naturelle du modèle continu). C'est la
valeur canonique injectée dans l'opérateur canalisé.
\end{proof}

\paragraph{Performance.}
Avec $\epsstar^{\mathrm{geom}}$ injecté sans ajustement, l'erreur médiane
sur $\sum_p \Ap\,p^{-s}$ est $0{,}085$ à $M = 1155$, $W = 24$ --- soit un
facteur $\sim 85$ au-dessus du baseline scalaire lissé adaptatif
($0{,}001$, qui utilise un $\beta$ tuné). L'apport n'est pas la précision
absolue mais l'\textbf{absence de paramètre libre}~: \eqref{eq:epsstar_geom}
remplace un tuning empirique par une covariance géométrique canonique.

\subsection{Formule de trace explicite}
\label{sec:trace_formula}

\begin{theoreme}[Formule de trace explicite]
\label{thm:trace_formula}
Pour tout $s$ avec $\Reop(s) > 1$,
\begin{equation}
\boxed{\;
-\,\dfrac{d}{ds} \log \Rres(s)
\;=\; \sum_{p} \dfrac{\kappap'(s)}{1 - \kappap(s)}
\;=\; \sum_{p} (\log p)\,\Big[\dfrac{1}{p^s - 1} \,+\, \Ap\,p^{-s}\Big]\,.\;}
\label{eq:trace_formula}
\end{equation}
\end{theoreme}

\begin{proof}
Par dérivation logarithmique de
$1 - \kappap(s) = (1 - p^{-s}) \exp(\Ap p^{-s})$. Posons $z = p^{-s}$,
de sorte que $dz/ds = -\log p \cdot z$. Un calcul direct donne
\[
\dfrac{\kappap'(s)}{1 - \kappap(s)}
\;=\; -(\log p)\,\Big[\Ap\,z - \dfrac{z}{1 - z}\Big]
\;=\; -(\log p)\,\Big[\Ap\,p^{-s} - \dfrac{1}{p^s - 1}\Big].
\]
En sommant sur $p$ et en changeant le signe global on obtient la formule
annoncée. Un \emph{sanity check} via $-\zeta'/\zeta(s) = \sum_n
\Lambda(n)\,n^{-s} = \sum_p (\log p)/(p^s - 1)$ donne la même expression
pour le premier terme.
\end{proof}

\paragraph{Vérification numérique.}
À $s = \sigma + 14.134725\,i$ avec $\sigma \in \{1.2,\,1.5,\,2,\,3\}$ et
$2\,000$ premiers tronqués, l'écart entre les deux formes
de~\eqref{eq:trace_formula} est de l'ordre de $10^{-26}$ (précision
\texttt{mpmath} 25 chiffres). L'identité est donc exacte au sens
algébrique.

\paragraph{Lecture Berry--Keating.}
L'opérateur de dilatation
$\HPTBK = (u\,p_u + p_u\,u)/2 = -i(u\partial_u + 1/2)$ engendre le flot
$e^{-it\HPTBK} f(u) = e^{-t/2}\,f(e^{-t}u)$. Sur la mesure atomique
$\nu = \sum_p \delta(u - \log p)$, les orbites périodiques du flot issu
de $\log p_0$ ont pour longueurs $\ell_p = k \log p_0$, $k \ge 1$,
correspondant aux puissances $p_0^{k}$. Le terme
$(\log p)/(p^s - 1) = \sum_{k \ge 1} (\log p)\,p^{-ks}$ s'identifie alors
à la somme sur ces orbites, et $\Ap\,p^{-s}$ représente la contribution
de couplage holonomique $\qplus/\qminus$ par premier.

\subsection{Trace régularisée opératorielle}
\label{sec:trace_reg}

L'écriture compacte de la formule de trace~\eqref{eq:trace_formula} en
termes de la résolvante de $\HPTBK$ et de l'opérateur $\DRop$ utilise une
projection de Mellin atomique
\begin{equation}
\Pi : L^{2}([\umin, \umax], du) \,\to\, \ell^{2}(\Primes_\gamma),
\qquad \Pi(f)_p \;=\; f(\log p)\,\sqrt{\log p},
\label{eq:Pi_def}
\end{equation}
où $\Primes_\gamma = \{p \text{ premier} : \umin \le \log p \le
\umax(\gamma)\}$. Le facteur $\sqrt{\log p}$ est la racine de la mesure
atomique de Haar multiplicative. La trace régularisée s'écrit alors
\begin{equation}
\boxed{\;
\Treg(s) \;:=\; \Tr_{\mathrm{reg}}\!\Big[(s - i\HPTBK)^{-1}\,\DRop(s)\Big]
\;=\; \sum_{p \in \Primes_\gamma} \kappap(s)\,(\log p)\,\Big[\dfrac{1}{p^s - 1} + \Ap\,p^{-s}\Big]\,.\;}
\label{eq:Treg}
\end{equation}
Dans $\Reop(s) > 1$, $\Treg(s) \equiv -d \log \Rres/ds$. Sur la ligne
critique $\Reop(s) = 1/2$, la somme prime n'est pas absolument
convergente ; elle est néanmoins distributionnellement convergente comme
élément de $\mathcal{S}'(\Reals_t)$, et la régularisation de Hadamard
donne un sens fonctionnel à $\Treg(1/2 + it)$.

\paragraph{Trois régularisations équivalentes.}
On utilisera trois régularisations interchangeables au sens distributionnel.
\begin{description}[nosep,leftmargin=*]
\item[($R_1$) Smearing par fonction test.] Pour $\phi_\eps$ une approximation
de l'unité de support compact à largeur $\eps$, on pose
$I_\eps(\gamma_0) := \int \phi_\eps(t - \gamma_0)\,F(1/2 + it)\,dt$.
Les pics $I_\eps(\gamma_n)$ se comportent comme $c/\eps$ quand $\eps \to 0^{+}$
(régularisation par lissage).
\item[($R_2$) Sommation Abel-Hadamard.] Avec poids $p^{-\eps}$, $\eps > 0$,
$F_\eps(t) := \sum_p (\log p)\,p^{-\eps}\,[1/(p^{1/2+it} - 1) + \Ap\,p^{-1/2-it}]$.
Par transformée de Fourier, $F_\eps$ est la convolution de $F$ par la
gaussienne $\phi_\eps$~: $R_2$ est équivalente à $R_1$ modulo le choix
de noyau régularisant.
\item[($R_3$) Soustraction du pôle principal.] $\widetilde F(s) := F(s) - 1/(s - 1)$
isole le pôle simple en $s = 1$ et applique l'extraction Hadamard à la
partie restante.
\end{description}

Les trois régularisations coïncident au sens distributionnel sur les
pôles de $\Rres$. La démonstration fonctionnelle stricte du prolongement
analytique de $\Treg$ à la ligne critique --- au-delà de la convergence
distributionnelle --- reste ouverte (\S\ref{sec:limites_ouvertures}).

\paragraph{Auto-correction : G\_HP4.d falsifié.}
Un premier test posait l'identité naïve « spectre direct de $\HPTBK$
borné par double barrière $= \{\gamma_n\}$ ». Ce test a été
falsifié~: le spectre direct est régulier (Bohr-Sommerfeld
asymptotiquement linéaire), alors que les $\gamma_n$ sont irréguliers, et
le ratio entre les deux spectres diverge. Cette falsification a motivé
le passage à la trace régularisée distributionnelle~: les zéros de
$\zeta$ sont des \emph{pôles distributionnels} de $\Treg(s)$, pas des
valeurs propres directes de $\HPTBK$. La discipline de falsifiabilité du
programme inclut cette correction au même titre que la falsification
Dirichlet de la \S\ref{sec:auto_correction_dirichlet}.

